%% Additional file 1 — Supplementary figures for the quantmsdiann MCP manuscript.
%% Build: cd paper && make supplementary

\documentclass[preprint,12pt]{elsarticle}

\usepackage[utf8]{inputenc}
\usepackage[T1]{fontenc}
\usepackage{microtype}
\usepackage{grffile}
\usepackage{amsmath}
\usepackage{url}
\usepackage{hyperref}
\usepackage{mcpfigures}
%% Shared macros for draft placeholders.
\newcommand{\placeholder}[1]{\textit{[#1]}}
\newcommand{\diannversion}[1]{\texttt{DIA-NN}~#1}
\newcommand{\quantmsdiann}{\texttt{quantmsdiann}}
\newcommand{\qpx}{\textsc{qpx}}


%% Number supplementary figures S1, S2, ... independently of the main text.
\renewcommand{\thefigure}{S\arabic{figure}}
\setcounter{figure}{0}

\journal{Molecular \& Cellular Proteomics}

\begin{document}

\begin{frontmatter}
\title{Additional file 1: Supplementary figures for ``quantmsdiann: a scalable SDRF-driven DIA-NN workflow for reanalysis of single-cell, spatial, and bulk proteomics datasets''}
\end{frontmatter}

\tableofcontents
\vspace{1em}

\section{Datasets benchmarked and reanalysed}
\label{sec:supp-datasets}

Table~\ref{tab:datasets-supp} lists every public dataset used in this study.
For each dataset, an SDRF was prepared from the original sample metadata and is provided alongside this supplementary deposition.
Raw files were downloaded from PRIDE and MassIVE using pridepy~\cite{kamatchinathan2025pridepy}.
For each cell-line reanalysis cohort, the original published quantification was downloaded from the PRIDE deposition and used as the comparator.
The cell-line bulk DIA cohorts were assembled by programmatic search of PRIDE Archive via the PRIDE API~\cite{deutsch2026proteomexchange}, filtering for projects annotated as DIA acquisition on human samples with at least one standardised cell line listed in the Cellosaurus reference, downloadable raw files, and SDRF-compatible sample metadata.

\begin{table}[htbp]
\centering
\small
\caption{Datasets benchmarked and reanalysed with \quantmsdiann{}.}
\label{tab:datasets-supp}
\begin{tabular}{llll}
\hline
\textbf{Accession} & \textbf{Instrument} & \textbf{Type} & \textbf{Public date} \\
\hline
\multicolumn{4}{l}{\textit{ProteoBench benchmarks}} \\
PXD049412 & Orbitrap Astral & Single-cell (Module~9) & 2024 \\
PXD062685 & timsTOF SCP & diaPASEF (Module~5) & 2025 \\
PXD070049 & ZenoTOF 7600 & Mixed-species (Module~10) & 2025 \\
Module 7  & Orbitrap Astral & Mixed-species (Module~7) & 2025 \\
\hline
\multicolumn{4}{l}{\textit{Bulk cell-line panel}} \\
PXD003539 & TripleTOF 5600 & Cell line (NCI-60) & 2016 \\
PXD030304 & TripleTOF 6600 & Cell line (ProCan) & 2022 \\
PXD004701 & TripleTOF 5600 & Cell line (Sun 2023) & 2016 \\
PXD041421 & timsTOF (diaPASEF) & Cell line (Wang 2023) & 2023 \\
PXD017199 & TripleTOF 6600 & Cell line (Tognetti 2021) & 2021 \\
\hline
\multicolumn{4}{l}{\textit{Scalability experiment}} \\
PXD071075 & Orbitrap Eclipse & Single-cell & 2025 \\
\hline
\end{tabular}
\end{table}

\section{Identification filtering and protein counting}
\label{sec:supp-filter}

All identification counts reported in the cell-line reanalyses and pan-cohort (main-text Figs.~3, 4 and Figs.~S5--S13 below) are target-only under a conservative contaminant/decoy filter: any row whose \texttt{Protein.Group} field contains a token prefixed with \texttt{CONTAM\_}, \texttt{Cont\_}, \texttt{ENTRAP\_}, or \texttt{DECOY\_} is dropped before counting, including mixed groups.
Original-paper headline counts were taken from the corresponding publication's table or text without re-filtering; where the original applied a per-protein peptide threshold (for example, the ProCan-DepMapSanger paper reports both an 8{,}498 protein-group headline at Global.Q.Value~$\leq 0.01$ and a 6{,}692 count under the additional $\geq 2$ supporting-peptide rule), the like-for-like \quantmsdiann{} count is reported under the same threshold.

\section{Per-process runtime and resource envelopes}
\label{sec:supp-performance}

These figures decompose the end-to-end wall-clock and resource use reported in main-text Fig.~1 by Nextflow process, derived from \texttt{trace.txt} records collected across the ProteoBench benchmark cohorts and the PXD071075 single-cell production sweep.

\begin{figure}[htbp]
  \centering
  \mcpincludegraphics[width=\mcponehalfcolwidth,keepaspectratio]{performance/runtime_per_step}
  \caption[Per-step runtime]{%
    \textbf{Per-step wall-clock distribution across Nextflow processes.}
    Distribution of task wall-clock time (boxplot, median and IQR) for each
    \quantmsdiann{} process across the ProteoBench benchmark cohorts. The
    per-file parallel stages (\texttt{FILE\_PREPARATION},
    \texttt{INSILICO\_LIBRARY\_GENERATION}, \texttt{PRELIMINARY\_ANALYSIS},
    \texttt{INDIVIDUAL\_ANALYSIS}) determine the workflow's
    parallelisation envelope, while the collective stages
    (\texttt{ASSEMBLE\_EMPIRICAL\_LIBRARY}, \texttt{FINAL\_QUANTIFICATION},
    \texttt{DIANN\_MSSTATS}) set the lower bound at maximum cluster width.
  }
  \label{fig:runtime-per-step}
\end{figure}

\begin{figure}[htbp]
  \centering
  \mcpincludegraphics[width=\textwidth,keepaspectratio]{performance/resources_per_step}
  \caption[Resources per task]{%
    \textbf{Peak memory and CPU envelope per Nextflow process.}
    Per-task peak resident-set-size memory (left panel) and CPU-hours
    (right panel) per \quantmsdiann{} process across the ProteoBench
    benchmark cohorts. The collective \texttt{FINAL\_QUANTIFICATION} stage
    dominates memory use on large cohorts; per-file stages remain well
    within the commodity 8--16~GB per-task envelope typical of HPC
    partitions.
  }
  \label{fig:resources-per-step}
\end{figure}

\section{ProteoBench head-to-head and accuracy benchmarks}
\label{sec:supp-proteobench}

These figures complement main-text Fig.~2 by placing \quantmsdiann{} on the public ProteoBench leaderboards alongside standalone \texttt{DIA-NN} GUI submissions and other DIA engines.

\begin{figure}[htbp]
  \centering
  \mcpincludegraphics[width=\textwidth,keepaspectratio]{quantmsdiann_benchmarks/main_diann_quantmsdiann_parity}
  \caption[ProteoBench parity]{%
    \textbf{ProteoBench parity between \quantmsdiann{} and the standalone
    \texttt{DIA-NN} GUI.}
    Head-to-head precursor and protein-group identification counts for the
    same raw data analysed by the official \texttt{DIA-NN} GUI submission to
    ProteoBench (grey) and by \quantmsdiann{} (blue). The two pipelines
    sit within $\pm$2\% on every module / version combination, confirming
    that the SDRF-driven \quantmsdiann{} wrapper does not erode \texttt{DIA-NN}
    performance.
  }
  \label{fig:proteobench-parity}
\end{figure}

\begin{figure}[htbp]
  \centering
  \mcpincludegraphics[width=\textwidth,height=\mcpmaxfigheight,keepaspectratio]{quantmsdiann_benchmarks/main_benchmarks_overview}
  \caption[ProteoBench per-version per-module breakdown]{%
    \textbf{ProteoBench identification counts per \texttt{DIA-NN} version
    and per module \textemdash{} detailed view.}
    Rows: ProteoBench DIA modules
    (Module~7 Orbitrap Astral 2~Th; Module~9 single-cell DIA, PXD049412;
    Module~5 timsTOF diaPASEF, PXD062685; Module~10 SCIEX ZenoTOF~7600,
    PXD070049).
    Columns: precursors (left, blue) and protein groups (right, red)
    quantified at 1\% FDR by \quantmsdiann{} when invoked with
    \diannversion{1.8.1}, \diannversion{2.1.0}, \diannversion{2.2.0},
    \diannversion{2.3.2}, and \diannversion{2.5.0} from the same SDRF.
    This is the per-version per-metric expansion of the headline summary
    in main-text Fig.~2a.
  }
  \label{fig:proteobench-overview}
\end{figure}

\section{Per-cohort completeness and overlap (cell-line reanalyses)}
\label{sec:supp-reanalyses}

These figures complement main-text Fig.~3 with per-condition completeness, per-protein peptide counts, gene/accession overlap, and per-cell-line missingness for each reanalysed cohort.

%% NOTE: The previous PXD003539 "genes per condition" panel sourced its
%% original-paper comparator from Walzer 2022 (E-PROT-73) rather than from
%% the Guo 2019 OpenSWATH analysis that the main-text PXD003539 reanalysis
%% uses, so the panel was not on the same comparator axis as main-text
%% Fig. 3a. The panel has been removed pending a Guo-2019-based replacement;
%% downstream Fig. S5 numbering has shifted accordingly.

\begin{figure}[htbp]
  \centering
  \mcpincludegraphics[width=\mcponehalfcolwidth,keepaspectratio]{PXD003539/supp_peptides_per_protein}
  \caption[PXD003539 peptides per protein]{%
    \textbf{PXD003539: distribution of peptides per protein group.}
    \quantmsdiann{} (DIA-NN, blue) versus Guo 2019 OpenSWATH (grey)
    on the same NCI-60 raw data.
  }
  \label{fig:supp-pxd003539-peptides}
\end{figure}

\begin{figure}[htbp]
  \centering
  \mcpincludegraphics[width=\mcponehalfcolwidth,keepaspectratio]{PXD003539/supp_venn_protein_accessions}
  \caption[PXD003539 accession overlap]{%
    \textbf{PXD003539: UniProt-accession overlap between original
    OpenSWATH analysis and \quantmsdiann{} reanalysis.}
  }
  \label{fig:supp-pxd003539-venn}
\end{figure}

\begin{figure}[htbp]
  \centering
  \mcpincludegraphics[width=\mcponehalfcolwidth,keepaspectratio]{PXD003539/supp_missing_values_per_run}
  \caption[PXD003539 missingness]{%
    \textbf{PXD003539: per-run missing-value fraction.}
    Distribution of per-run missingness across the NCI-60 panel for both
    pipelines.
  }
  \label{fig:supp-pxd003539-missing}
\end{figure}

\begin{figure}[htbp]
  \centering
  \mcpincludegraphics[width=\mcponehalfcolwidth,keepaspectratio]{PXD004701/supp_proteins_per_subtype}
  \caption[PXD004701 per subtype]{%
    \textbf{PXD004701 (Sun 2023): protein groups per breast-cancer subtype.}
    Distribution of identifications across the 76 cell lines stratified by
    TNBC / non-TNBC subtype.
  }
  \label{fig:supp-pxd004701-subtype}
\end{figure}

\begin{figure}[htbp]
  \centering
  \mcpincludegraphics[width=\mcponehalfcolwidth,keepaspectratio]{PXD004701/supp_missing_values_per_run}
  \caption[PXD004701 missingness]{%
    \textbf{PXD004701 (Sun 2023): per-run missing-value fraction.}
  }
  \label{fig:supp-pxd004701-missing}
\end{figure}

\begin{figure}[htbp]
  \centering
  \mcpincludegraphics[width=\mcponehalfcolwidth,keepaspectratio]{PXD004701/main_comparison}
  \caption[PXD004701 reanalysis vs original]{%
    \textbf{PXD004701 (Sun 2023, PCT-SWATH breast-cancer panel, 76 lines):
    \quantmsdiann{} reanalysis versus the originally published quantification.}
    Under the original consistency filter ($-90\%$ missing, proteotypic)
    \quantmsdiann{} recovers 6{,}146 versus 6{,}091 protein groups
    ($+0.9\%$); raw-matrix counts are within 10\% on both
    proteotypic-peptide and protein-group metrics. This near-parity cohort
    --- already analysed with a current DIA engine in the original study ---
    is shown here rather than in the main reanalysis figure.
  }
  \label{fig:supp-pxd004701-comparison}
\end{figure}

\begin{figure}[htbp]
  \centering
  \mcpincludegraphics[width=\textwidth,keepaspectratio]{plexDIA/MSV000093870/main_galatidou_comparison}
  \caption[MSV000093870 single-cell plexDIA reanalysis]{%
    \textbf{MSV000093870 (Galatidou 2024, single-cell oocyte plexDIA):
    \quantmsdiann{} reanalysis versus the original.}
    (\textbf{a})~Protein groups per single cell (Galatidou $n=97$ post-QC
    oocytes; \quantmsdiann{} $n=107$ channel-confident cells): comparable
    median depth (1{,}784 versus 1{,}736). (\textbf{b})~Total protein groups
    across the dataset: 2{,}122 (original) versus 2{,}904 (\quantmsdiann{}).
    (\textbf{c})~UniProt-accession overlap: 2{,}013 shared, 109 original-only,
    1{,}539 \quantmsdiann{}-only --- 95\% of the original recovered plus
    1{,}539 additional proteins from the same raw data. This figure
    underpins the single-cell panel of the main reanalysis figure.
  }
  \label{fig:supp-msv093870-plexdia}
\end{figure}

\begin{figure}[htbp]
  \centering
  \mcpincludegraphics[width=\mcponehalfcolwidth,keepaspectratio]{PXD030304/supp_proteins_per_tissue}
  \caption[PXD030304 per tissue]{%
    \textbf{PXD030304 (ProCan-DepMapSanger): protein groups per tissue.}
    Distribution of identifications across the 949 cell lines stratified by
    tissue of origin.
  }
  \label{fig:supp-pxd030304-tissue}
\end{figure}

\begin{figure}[htbp]
  \centering
  \mcpincludegraphics[width=\mcponehalfcolwidth,keepaspectratio]{PXD030304/supp_missing_values_per_run}
  \caption[PXD030304 missingness]{%
    \textbf{PXD030304: per-run missing-value fraction across the
    949-cell-line panel.}
  }
  \label{fig:supp-pxd030304-missing}
\end{figure}

\bibliographystyle{elsarticle-num}
\bibliography{references}

\end{document}
